\chapter{Building, packaging, shipping}
\label{ch:build}

This last chapter is the practical mile: how to compile \maya{}, how
to depend on it from your own project, how to ship a binary that
runs anywhere, and what to do when something goes wrong. It assumes
you have read the architecture chapters and want to put a real
program on a real machine.

\section{Compiler requirements}

\maya{} requires \cpp{}26. As of this writing the meaningful
toolchains are:

\begin{center}
\small
\begin{tabular}{lll}
\textbf{Compiler} & \textbf{Minimum} & \textbf{Notes} \\
\hline
GCC                & 15 & recommended on Linux and macOS \\
Clang              & 18 & needs libstdc++ 15 or libc++ trunk \\
MSVC               & 14.40 & \cpp{}23 + \texttt{/std:c++latest} \\
AppleClang         & --- & not yet \cpp{}26-capable; use Homebrew GCC \\
\end{tabular}
\end{center}

On Apple Silicon this means \code{brew install gcc} then explicitly
configure with \code{cmake -DCMAKE\_CXX\_COMPILER=g++-15}. The
AppleClang shipped with Xcode is not ready for \cpp{}26 yet.

\section{Building \maya{}}

The standard incantation:

\begin{lstlisting}[style=shell]
git clone --recursive https://github.com/1ay1/agentty
cd agentty/maya
cmake -B build
cmake --build build -j$(nproc)
\end{lstlisting}

A few useful CMake options:

\begin{itemize}[leftmargin=*]
  \item \code{-DMAYA\_BUILD\_EXAMPLES=ON} --- build the 26 example
    programs. On by default if \maya{} is the top-level project; off
    when it is a sub-build.
  \item \code{-DMAYA\_BUILD\_TESTS=ON} --- build the test suite.
    Run with \code{ctest --test-dir build}.
  \item \code{-DCMAKE\_BUILD\_TYPE=Release} --- enable
    \code{-O3 -DNDEBUG}. \maya{} also turns on link-time
    optimisation (LTO) automatically in Release.
  \item \code{-DCMAKE\_INTERPROCEDURAL\_OPTIMIZATION=OFF} ---
    disable LTO if your linker is short on memory.
\end{itemize}

A clean Release build of \maya{} plus all examples takes about a
minute on a modern laptop.

\section{Linking \maya{} into your own program}

The supported path is to install \maya{} and \code{find\_package} it
from your project's CMake.

\subsection{Installing}

\begin{lstlisting}[style=shell]
cmake -B build -DCMAKE_BUILD_TYPE=Release
cmake --build build -j
cmake --install build --prefix /usr/local
\end{lstlisting}

This installs headers under \code{\$prefix/include/maya/} and the
static library under \code{\$prefix/lib/}, plus a CMake config
package under \code{\$prefix/lib/cmake/maya/}.

\subsection{Depending on it}

In your project's \code{CMakeLists.txt}:

\begin{lstlisting}[style=shell]
cmake_minimum_required(VERSION 3.28)
project(my_app LANGUAGES CXX)

set(CMAKE_CXX_STANDARD 26)
set(CMAKE_CXX_STANDARD_REQUIRED ON)

find_package(maya 0.1 REQUIRED)

add_executable(my_app main.cpp)
target_link_libraries(my_app PRIVATE maya::maya)
\end{lstlisting}

Then in \code{main.cpp}:

\begin{lstlisting}
#include <maya/maya.hpp>
#include <maya/widget/markdown.hpp>     // each widget header opt-in

int main() {
    maya::print(maya::dsl::t<"Hello from your app">.build());
}
\end{lstlisting}

\subsection{Vendoring via \code{FetchContent}}

If you prefer to vendor \maya{} alongside your sources:

\begin{lstlisting}[style=shell]
include(FetchContent)
FetchContent_Declare(maya
    GIT_REPOSITORY https://github.com/1ay1/agentty.git
    GIT_TAG        master
    SOURCE_SUBDIR  maya
)
set(MAYA_BUILD_EXAMPLES OFF CACHE BOOL "" FORCE)
set(MAYA_BUILD_TESTS    OFF CACHE BOOL "" FORCE)
FetchContent_MakeAvailable(maya)

target_link_libraries(my_app PRIVATE maya::maya)
\end{lstlisting}

This downloads, builds, and links \maya{} as part of your
configure step. No installation required. Slower first build,
faster CI.

\section{Sandboxing your program (Linux / macOS)}

\maya{} itself does not sandbox; that is a feature of programs like
\agentty{} that wrap shell calls. If your TUI shells out, look at
\agentty{}'s \file{src/tool/util/sandbox.cpp} as a reference. The
relevant building blocks:

\begin{itemize}[leftmargin=*]
  \item \textbf{Linux:} \code{bwrap} (bubblewrap) with read-only
    binds for \code{/usr}, \code{/etc}, etc., a fresh
    \code{tmpfs} for \code{/tmp}, a shared network namespace.
  \item \textbf{macOS:} \code{sandbox-exec} with a sandbox profile
    file written to a temp path and removed on exit.
\end{itemize}

The pattern is to detect the backend by attempting a
\code{<exe> --version} probe, fall back to no-sandbox if absent,
and surface the decision to the user.

\section{Producing a static binary}

If you want one binary you can scp to any compatible machine,
build with static linking. The shape (Linux) is:

\begin{lstlisting}[style=shell]
cmake -B build-static \
    -DCMAKE_BUILD_TYPE=Release \
    -DCMAKE_EXE_LINKER_FLAGS="-static-libstdc++ -static-libgcc"
cmake --build build-static -j
\end{lstlisting}

This statically links libstdc++ and libgcc, but leaves libc
dynamic (fully-static glibc breaks DNS resolution via NSS). For a
100\%-static binary on Linux you need a musl toolchain
(\code{x86\_64-linux-musl-g++} from Alpine or
\code{musl-cross-make}). On macOS, only \code{libSystem} is
dynamic --- that is Apple's only supported ABI.

On Windows, link the static MSVC runtime (\code{/MT}) and use the
\code{x64-windows-static} vcpkg triplet for any third-party
dependencies. \agentty{}'s release pipeline does exactly this; you
can crib from its \code{CMakeLists.txt}.

\section{Packaging}

For distribution, see \agentty{}'s \file{packaging/} directory and
\file{scripts/release.sh}. The script builds:

\begin{itemize}[leftmargin=*]
  \item A \code{.deb} (Debian / Ubuntu).
  \item An \code{.rpm} (Fedora / RHEL).
  \item A \code{.pkg.tar.zst} (Arch).
  \item A Homebrew formula.
  \item A Scoop manifest (Windows).
  \item An AUR PKGBUILD.
  \item Raw static binaries.
\end{itemize}

The whole pipeline runs locally; no CI required. The script is a
worked example for any project that wants a similar one-command
release flow.

\section{Debugging}

\subsection{Terminal stuck in raw mode}

If your program crashed and the shell is acting weird, run
\code{stty sane} or \code{reset}. Avoid this in the future by using
RAII for the terminal state (which \maya{} does internally).

\subsection{Rendering looks corrupted}

First, check that your terminal supports truecolour and the modes
\maya{} uses. \code{echo \$TERM} should show a modern value
(\code{xterm-256color}, \code{kitty}, \code{wezterm}). Some old
terminals (Windows command prompt) need explicit enablement of
ANSI processing.

Second, capture a frame to a file:

\begin{lstlisting}[style=shell]
script -q -c './my_program' frames.txt
\end{lstlisting}

The \code{frames.txt} file contains every byte your program sent;
inspect it with \code{cat -v frames.txt} to see the raw ANSI.

\subsection{Slow rendering}

If frames are dropping, the bottleneck is almost certainly your
view function or your event handler, not \maya{}. Profile with
\code{perf record} (Linux), \code{Instruments} (macOS), or
\code{Tracy} (cross-platform). \maya{}'s render path itself
measures in the tens of microseconds.

\section{Going further}

This book has been a guided tour. The repository contains a great
deal more:

\begin{itemize}[leftmargin=*]
  \item \file{maya/docs/}: 14 markdown chapters that overlap with
    this book but are more reference-style.
  \item \file{maya/examples/}: 26 fully-worked programs ranging
    from \file{counter.cpp} (60 lines) to \file{messenger.cpp}
    (3000+ lines, a real chat UI).
  \item \file{maya/docs/internals/}: deep engineering writeups on
    the corner cases of inline rendering, the witness chain, and
    scrollback integrity.
  \item \agentty{}'s repository: the largest production user of
    \maya{}, with its own design discipline that complements the
    library's.
\end{itemize}

If you finished this book and built something with \maya{}, you have
joined a small club of people who write terminal UIs by composing
typed values rather than by typing escape sequences. That is the
right way to build them; the rest of the world will get there
eventually.

\begin{recap}
Build \maya{} with CMake against GCC 15+ or Clang 18+ (or MSVC
14.40+ on Windows). Use \code{find\_package(maya)} or
\code{FetchContent} to depend on it. Static binaries need
\code{-static-libstdc++ -static-libgcc} on Linux (musl for fully
static) or \code{/MT} on Windows. Packaging is one script away in
\agentty{}'s \file{scripts/release.sh}. When something looks wrong,
\code{script(1)} the byte stream and read the ANSI.
\end{recap}

\section*{Workshop}

\begin{enumerate}[leftmargin=*]
  \item Set up a new project, vendor \maya{} via
    \code{FetchContent}, and write a one-screen TUI for something
    you actually want --- a directory size visualiser, a small
    log tail, an HTTP-endpoint poller.
  \item Build it static. Scp it to a different machine. Run it.
    Note how the dependency story stays exactly: zero.
  \item Contribute back. Pick a widget that does not yet exist
    (a kanban board, a piano-roll, a syntax-highlighted JSON
    viewer) and PR it. The maintainers are friendly; the
    architecture makes the work tractable.
\end{enumerate}
